\section*{Discussion}

We set out to test a claim that the risk-sensitive control literature has largely assumed: that shaping the tail of a learned policy's outcome distribution makes an autonomous delivery operation more resilient when conditions drift away from those it was trained in. In a multi-UAV parcel delivery setting with wind, degraded navigation, mixed cooperative and non-cooperative traffic and uncertain energy, the claim did not hold. Three tail-risk designs, a fixed conditional value-at-risk objective of the kind used by risk-conditioned and worst-case actor-critics, the same objective with a hazard-adaptive risk level, and a variant that places the tail measure on a learned catastrophe-cost distribution rather than on the return, all finished within the bootstrap and seed-to-seed spread of a well-tuned risk-neutral soft actor-critic, ahead of it in some conditions and behind it in others. Turning the deployment risk level towards caution did not lower collision rates, and the ablations show that neither the choice of hazard cues nor the placement of the risk measure changes that conclusion.

Three mechanisms explain the result, and each of them is visible in the data rather than inferred. First, the lower tail of the return distribution in this task is not made of collisions. Delivery returns vary because wind, payload and traffic change flight time and energy, and this ordinary variability occupies the same quantiles that a moderate risk level examines, so an agent that maximises a tail statistic mostly buys insurance against being slow. Placing the measure on a catastrophe-cost distribution removes that confusion in principle, and the risk level then acts as a multiplier on the probability of a crash, but it introduces a second problem: the cost signal is sparse and its tail must be estimated in exactly those states where data are scarcest. Second, caution expresses itself as hesitancy. A risk-averse policy slows down and holds back near conflicts, which lengthens the time spent in proximity to other aircraft rather than increasing the geometric margin, and exposure and margin trade against each other almost exactly in our measurements. Third, and most important, a large part of the residual risk is not controllable by the policy at all. Under the sensor shift one intruder in five is undetected at any instant, and no property of the agent's own return distribution carries information about an aircraft it cannot see. Tail-risk learning can only reallocate behaviour among observed contingencies.

This does not contradict the studies that report gains from risk-sensitive navigation, but it does bound their scope. Those results come predominantly from single-agent settings with dense static clutter, full observability of the obstacles that matter, no energy feasibility constraint and evaluation inside the training distribution; several of them compare against risk-neutral variants of the same distributional learner rather than against a strong off-policy baseline trained to the same budget. Our environment differs in all four respects. The lesson we draw is methodological: risk sensitivity is a property of an objective, not of a system, and a claim of resilience should be tested against a strong risk-neutral baseline, on paired scenarios, with shift applied one mechanism at a time.

Where the failure mechanism is understood and a conservative fallback exists, statistical calibration was decisive. The return-to-home decision has both properties: the energy required to reach a pad is a modelled quantity with a monotone dependence on a scalar reserve, and aborting a delivery is always available. The online conformal reserve held its specified coverage of return-leg energy in every shift family, while fixed margins lost a third of their coverage as conditions moved, and it removed energy depletion almost entirely under the energy shift. The guarantee behind this behaviour is deliberately weak in what it assumes and strong in what it delivers: Proposition~\ref{prop:coverage} requires no exchangeability, no model of the shift and no probabilistic assumption of any kind, and it holds for the leg-level quantity the operator actually cares about. The fleet-level calibrator matters here, because it pools evidence across aircraft and adapts within a few hundred legs, which is minutes of operation for a fleet of realistic size.

Calibration is not free, and we report where it loses. The reserve buys safety with abandoned deliveries, and the exchange rate depends on the condition. Under the energy shift a conformal target of five percent miscoverage dominates the fixed margin that reaches comparable safety, delivering more parcels at lower depletion, and it does so simultaneously at nominal conditions, which no single fixed margin manages. Under the wind shift, however, a moderate fixed margin is at least as good, because strong wind inflates the modelled return energy itself rather than the discrepancy between model and measurement that the calibrator tracks. An operator should therefore treat the miscoverage target as a policy variable, chosen from a frontier that our sweep makes explicit, and should expect the benefit of calibration to be largest when the model error, rather than the model input, is what drifts.

Two practical directions follow for the failure mode that neither mechanism addressed. The first is detection and cooperation rather than objective shaping: most residual losses in our study are conflicts with aircraft that are seen late or not at all, so calibrated prediction of other aircraft's motion coupled to an explicit resolution rule, or a cooperative separation protocol of the kind envisaged in traffic-management concepts, attacks the problem where it lives. The second is the reserve idea applied to separation: a conformal bound on the miss distance, with a backup manoeuvre as the fallback, has the same structure that worked for energy, a modelled quantity, a monotone knob and an always-available conservative action. We would expect that combination to move collision rates in a way that tail-risk objectives did not.

Finally, a note on how this study was conducted. Our first tail-risk design placed the risk measure on the return distribution, and when it failed to improve on the baselines we designed the catastrophe-cost variant in response. That second design was fixed before it was evaluated, its weight was set a priori from a scale argument rather than tuned, the deployment risk level was selected on nominal validation districts, and no test condition was used for any design decision; both designs are reported. We record the sequence because a study whose conclusion is that a popular mechanism does not help should be transparent about how hard it tried to make it work.
